% !TEX root = 00_instructivo.tex
\chapter{Sensores de estímulos eléctricos}
\section{\obj}
\capacidad
\begin{itemize}
    \item Programar un Arduino con funciones básicas como introducción al microcontrolador.
    \item Calcular el error asociado a la medición utilizando una multimetro digital de alta precisión.
\end{itemize}

\section{\mat}
\textbf{A suministrar por la Escuela:}
\begin{itemize}
    \item 1 ARDUINO UNO R4 MINIMA
    \item Cable USB de conexión del Arduino
    \item 1 potenciómetro (5 k$\Omega$ o más)
    \item 2 LEDs
    \item 2 resistencias de 330 $\Omega$ (o un valor cercano)
    \item Cables de interconexión (jumpers) macho-macho
\end{itemize}

\section{\pro}
\begin{enumerate}
        
    \item Realice las conexiones del circuito tal y como se indica en la Figura \ref{fig:elec1}.
    \item Abra el programa Arduino IDE
    
    % Requiere en el preámbulo:
    % \usepackage{circuitikz}
    % \usepackage{float}
    % \usepackage{siunitx}

    % Requiere en el preámbulo:
    % \usepackage{circuitikz}
    % \usepackage{float}
    % \usepackage{siunitx}

    \begin{figure}[H]
    \centering
    \resizebox{1\textwidth}{!}{
    \begin{circuitikz}
    \tikzstyle{every node}=[font=\LARGE]
    \draw (8.25,14.75) to[short, -o] (8.25,14.75) ;
    \draw  (1.5,16.5) rectangle (7.5,6.25);
    \draw (7,14.75) to[short, -o] (7,14.75) ;
    \draw (7.25,15.25) to[short, -o] (8,15.25) ;
    \draw (7.25,11.5) to[short, -o] (8,11.5) ;
    \draw (7.25,8) to[short, -o] (8,8) ;
    \draw (1.75,15) to[short, -o] (1,15) ;
    \draw (1.75,8.5) to[short, -o] (1,8.5) ;
    \draw (1.75,12) to[short, -o] (1,12) ;
    \draw (-2.25,9.75) to[R,l={ \LARGE 5 k$\Omega$}] (-2.25,13.5);
    \draw [-{Stealth}] (-3,11) -- (-1.5,12);
    \node [font=\normalsize] at (2.25,15) {5 V};
    \node [font=\normalsize] at (6.75,15.25) {12};
    \node [font=\LARGE] at (4.75,11.75) {Arduino};
    \node [font=\normalsize] at (6.75,11.5) {3};
    \node [font=\normalsize] at (6.75,8) {GND};
    \node [font=\normalsize] at (2.25,8.5) {GND};
    \node [font=\normalsize] at (2.25,12) {A0};
    \draw (-2.25,13.5) to[short] (-2.25,15);
    \draw (-2.25,15) to[short] (1,15);
    \draw (-1.5,12) to[short] (1,12);
    \draw (-2.25,9.75) to[short] (-2.25,8.5);
    \draw (1,8.5) to[short] (-2.25,8.5);
    \draw (9,11.5) to[R,l={ \normalsize 330 $\Omega$}] (11,11.5);
    \draw (9.25,15.25) to[R,l={ \normalsize 330 $\Omega$}] (11.5,15.25);
    \draw (12,10.5) to[empty led] (12,8);
    \draw (14.5,11.5) to[empty led] (14.5,8);
    \draw (8,8) to[short] (14.5,8);
    \draw (8,15.25) to[short] (9.25,15.25);
    \draw (14.5,11.5) to[short] (14.5,15.25);
    \draw (11.5,15.25) to[short] (14.5,15.25);
    \draw (8,11.5) to[short] (9,11.5);
    \draw (12,10.5) to[short] (12,11.5);
    \draw (11,11.5) to[short] (12,11.5);
    \end{circuitikz}
    }
    \caption{Circuito de Arduino con potenciómetro y dos LEDs.}
    \label{fig:elec1}
    \end{figure}

    \item Añada en la función \textit{loop()} del código una sección para leer el valor analógico del pin A0 e imprimirlo a la consola serial. Córralo y verifique que al variar el potenciómetro, varíe correctamente el valor.
    \item Modifique el \textit{loop()} y agregue un delay, que varíe linealmente entre 100 ms y 1 s, dependiendo de la posición del potenciómetro. Córralo y verifique que ahora las lecturas se impriman en la terminal con diferentes tiempos de retraso, variando con el valor del potenciómetro.
    \item Agregue a la función \textit{loop()} el código necesario para que el LED en el pin 12 se encienda y se apague con el retraso dado (para esto hay que agregar un segundo delay para controlar el tiempo de apagado). Córralo y verifique que los tiempos de encendido y apagado del LED varían correctamente al variar el potenciómetro.
    \item Agregue el código necesario en \textit{loop()} para que el brillo del LED varíe también con el valor leído del potenciómetro. Córralo y verifique que el brillo del LED varíe correctamente al cambiar el valor del potenciómetro.
    \item Mida la tensión entre la tierra y el pin A0 (el pin central del potenciómetro) con un multímetro y compare las lecturas obtenidas con el Arduino y el multímetro (al menos 5 veces).
    % {\scriptsize 
    %     \lstinputlisting[language=Arduino, numbers=none, showstringspaces=false]{L01/L01.ino}
    % }


    \section{Preguntas para el analisis de resultados}
    \begin{itemize}
        \item ¿Cuál es el rango de valores que devuelve \textit{analogRead()} en un Arduino UNO?
        \item Con el resultado anterior, ¿cuántos bits de resolución tiene el ADC del Arduino UNO?
        \item ¿Cómo se convierte de los valores obtenidos con \textit{analogRead()} a una medición de tensión eléctirca?
        \item ¿Cómo calculó la fórmula para ajustar los valores obtenidos de \textit{analogRead()} a valores entre 100 ms y 1 s para la función \textit{delay()}?
        \item Calcule el error relativo y el error absoluto para cada una de las 5 lecturas de tensión entre el Arduino y el potenciómetro (asuma que el multímetro es el valor real).
        \item Incluya el código que realizó en esta práctica en el informe.
    \end{itemize}
\end{enumerate}
